%*****************************************************************************
% Layout-Vorschlag zur Erstellung von Bachelor-, Master oder Diplomarbeiten  *
% Modifikation moeglich                                                      *
% S. Gippert, 2012                                                           *
% ****************************************************************************

\documentclass[11pt,a4paper]{article}
\usepackage[T1]{fontenc}   
\usepackage[utf8]{inputenc}
\usepackage{lmodern}
\usepackage{algorithm2e}
\usepackage{graphicx}
\usepackage[T1]{fontenc}   
\usepackage[ngerman]{babel}
\usepackage{float}
\usepackage{amsmath}
\usepackage{amssymb}
\usepackage{caption}
\usepackage{mathalpha}
\usepackage{amsthm}
\usepackage[usenames,dvipsnames]{xcolor}
\usepackage{multirow}
\usepackage{longtable}
\usepackage{tabularx}
\usepackage{setspace}
%\usepackage{listings}
\usepackage{listingsutf8}
\usepackage{adjustbox}
\usepackage{geometry}
 \geometry{
 a4paper,
 bindingoffset=6mm,
 total={160mm,240mm},
 left=25mm,
 top=25mm,
 }
\usepackage{parskip}
\usepackage[backend=biber,style=numeric
]{biblatex}
\addbibresource{thesisbib.bib}
\usepackage{pdflscape}
\usepackage{url}
\usepackage{tikz}
\usepackage{pgfplots}
\usepackage{pgfplotstable}
\pgfplotsset{compat=1.18}
\usepackage[simplified]{pgf-umlcd}
\usetikzlibrary{shapes.geometric}
\usetikzlibrary{positioning}
\usepackage{tocloft}
\usepackage[colorlinks=true, linkcolor=black, urlcolor=black, citecolor=black]{hyperref}

\definecolor{c1}{RGB}{31,119,180} 
\definecolor{c2}{RGB}{255,127,14}
\definecolor{c3}{RGB}{44,160,44}  
\definecolor{c4}{RGB}{214,39,40}
\definecolor{c5}{RGB}{148,103,189} 
\definecolor{c6}{RGB}{23,190,207}
\definecolor{c7}{RGB}{140,86,75}  
\definecolor{c8}{RGB}{227,119,194}
\definecolor{c9}{RGB}{188,189,34}

\DeclareMathAlphabet{\mathdutch}{U}{dutchcal}{m}{n}

\renewcommand{\umltextcolor}{black} % Textfarbe
\renewcommand{\umldrawcolor}{black}  % Rahmen- und Linienfarbe
\renewcommand{\umlfillcolor}{white}

\DefineBibliographyStrings{ngerman}{
  andothers = {et\addabbrvspace al\adddot},
  andmore   = {et\addabbrvspace al\adddot},
}
\newcommand{\umlagg}{\tikz[baseline=-0.5ex]{
    \draw (0,0) -- (0.4,0); 
    \draw (0.4,0) node[draw, diamond, inner sep=1.5pt, anchor=west] {};
}}

% Befehl für den gestrichelten Pfeil
\newcommand{\umldash}{\tikz[baseline=-0.5ex]{
    \draw[dashed, ->] (0,0) -- (0.6,0);
}}

\widowpenalty=10000
\clubpenalty=10000 

\newenvironment{Beweis}{\begin{proof}[Beweis]}{\end{proof}}
%\theoremstyle{captioned}
\newtheorem{Definition}{Definition}[section]% Definitionen
\newtheorem{Satz}{Satz}[section]% Definitionen
\def\Real{{{\rm I\!R}}}

\DeclareMathOperator{\dist}{dist}
\DeclareMathOperator{\midi}{mid}
\DeclareMathOperator{\rank}{Rang}
\DeclareMathOperator{\integer}{integer\,}
\DeclareMathOperator{\subjectto}{s.t.\,}
\DeclareMathOperator{\argmin}{argmin\,}
\DeclareMathOperator{\erf}{erf\,}
\DeclareMathOperator{\e}{e}
\newcommand{\sigmoid}{\zeta}
\newcommand{\E}[1]{\mathbb E \left(#1\right)}
\newcommand{\Var}[1]{\textup{Var} \left(#1\right)}
\newcommand{\pdf}{Wahrscheinlichkeitsdichtefunktion }
\newcommand{\cdf}{Verteilungsfunktion }
\newcommand{\Prob}[1]{\mathbb P\left(#1\right)}

\setlength{\tabcolsep}{3pt}

\newcounter{rowcount}
\newcommand{\nextrow}{%
    \stepcounter{rowcount}%
    \ifnum\value{rowcount}=5
        \setcounter{rowcount}{0}%
        \\ \hline % Nach 5 Zeilen: Umbruch erlaubt + Linie
    \else
        \\* % Sonst: Umbruch verboten
    \fi
}

% Dokumentbeginn
\begin{document}

% Titelseite ----------ANPASSEN!----------------------
\begin{titlepage}
\title{{\textbf{Report}} \\[1.5cm] \textbf{Project 1: Jacobi Sweeps with CUDA} \\[1cm]}
\author{\textbf{Group 2} \\Laura Arango \\Andreas Serin}
\date{\today}
\maketitle
\vfill 
% Weitere Angaben (Arbeitsgruppe, Institut, Universität, Betreuung usw.)
\begin{center}
\begin{figure}[h]
\centering
\includegraphics[width=4.75cm]{siegel}
%\includegraphics[scale=0.25]{uzk_logo} % bunt
\end{figure} 
\begin{large}
%Lehrstuhl für Visualisierung und Visual Analytics \\[3pt]
Department Mathematik/Informatik \\[2pt]
Mathematisch-Naturwissenschaftliche Fakultät \\[2pt]
Universität zu Köln \\[1.5cm]
Betreuer: Dr. Martin Lanser \\[2pt]
%Erstgutachter: Dr. Alexander Apke
\end{large}
\end{center}
\vfill 
% Unterdrücken der Seitennummerierung
\thispagestyle{empty}
\end{titlepage}
\clearpage \thispagestyle{empty} \cleardoublepage 

\makeatletter
\def\l@figure#1#2{\@dottedtocline{1}{0em}{4ex}{#1}{#2}}
\def\l@table#1#2{\@dottedtocline{1}{0em}{4ex}{#1}{#2}}
 \def\l@lstlisting#1#2{\@dottedtocline{1}{0em}{4ex}{#1}{#2}}
  \def\l@algocf#1#2{\@dottedtocline{1}{0em}{4ex}{#1}{#2}}
\makeatother

\singlespacing

\setlength{\cftsecindent}{0.75em}    % Vorher ca. 1.5em
\setlength{\cftsubsecindent}{1.9em}  % Vorher ca. 3.8em
\setlength{\cftsubsubsecindent}{3.5em} % Vorher ca. 7.0em

\pagestyle{plain}
\cleardoublepage
\setcounter{page}{1}
%\pagenumbering{roman}
%\thispagestyle{plain}
%\tableofcontents
%\newpage
%\thispagestyle{plain}

%\cleardoublepage
%\addcontentsline{toc}{section}{\listfigurename}
%\listoffigures
%\newpage
%\cleardoublepage
%\addcontentsline{toc}{section}{\listtablename}
%\listoftables
%\newpage

\renewcommand{\listalgorithmcfname}{Verzeichnis der Algorithmen}%
\renewcommand{\algorithmcfname}{Algorithmus}%
\renewcommand{\algorithmautorefname}{\algorithmcfname}%

%\renewcommand{\lstlistingname}{Datei}
%\renewcommand{\lstlistlistingname}{Verzeichnis der Dateien}


%\renewcommand\cftlistalgorithmsaftersnum{:}
%\renewcommand\cftalgorithmscfpresnum{Algorithmus~}
%\cftsetindents{algorithms}{1.5em}{7em}


%renewcommand\cftlstlistlistingaftersnum{:}
%renewcommand\cftlstlistlistingpresnum{Datei~}
%\cftsetindents{lstlistlisting}{1.5em}{7em}



%\cleardoublepage
%\addcontentsline{toc}{section}{\listalgorithmcfname} 
%\listofalgorithms
%\newpage


%\addcontentsline{toc}{section}{\lstlistlistingname}
%\lstlistoflistings

%\newpage
%\cleardoublepage
%\thispagestyle{plain}
%\addcontentsline{toc}{section}{Abkürzungsverzeichnis}
\section*{Abkürzungsverzeichnis}
\begin{tabbing}
xxxxxxxxxxxx\=xxxxxxxxxxxxxxxxxxxxxxxxxxxxxxx\kill
BET \>\texttt{BETWEENNESS} \\
BFI \>Brute-Force-Iteration \\
CLO \>\texttt{CLOSENESS} \\
CSP \>Constraint Satisfaction Problem \\
DEG \>\texttt{DEGREE} \\
EIG \>\texttt{EIGENVECTOR} \\
GBP \>Graph Burning Problem \\
HAR \>\texttt{HARMONIC} \\
ILP-PROP \>Integer Linear Programming based on propagation \\
KAD \>\texttt{KADABRA} \\
KAT \>\texttt{KATZ} \\
PAG \>\texttt{PAGERANK} \\
RAN \>\texttt{RANDOM} \\
SSSP \>Single-Source Shortest Path \\
%SQP \>Sequential Quadratic Programming \\
\end{tabbing}

 
%\newpage
%\cleardoublepage
%\thispagestyle{plain}
%\addcontentsline{toc}{section}{Notation}
\section*{Notation}
\subsection*{Allgemein}
\begin{tabbing}
xxxxxxxxxxxxxxxxxx\=xxxxxxxxxxxxxxxxxxxxxxxxxxxxxxxxxxxxx\kill
$t$, $u$, $x$, $z$, $\theta$, $\theta_0$, $\theta_1$ \>Reelle Parameter \\
$i, j, k, l, m, n$ \>Natürliche Zahlen \\
$\arctan$ \>Arkustangens \\
$\erf$ \>Gaußsche Fehlerfunktion \\
$g(x)$, $\phi(x)$ \>Funktionen \\
\end{tabbing}
\subsection*{Graphen}
\begin{tabbing}
xxxxxxxxxxxxxxxxxx\=xxxxxxxxxxxxxxxxxxxxxxxxxxxxxxxxxxxxx\kill
$b(G)$ \>Brennzahl von $G$ \\
$b_i$ \>$i$-ter Knoten einer Brennsequenz \\
$b_{i,j}, s_{i,j}$, $x_j$ \>Binärvariablen der ILP-PROP-Formulierung \\
$D$ \>Durchmesser des Graphen $G$ \\
$E$ \>Kantenmenge \\
$G = (V, E)$ \>Graph \\
$H_s(v)$ \>Höhe des Knotens $v$ im BFS-Baum mit Startknoten $s$ \\
$K_n$ \>Vollständiger Graph mit $n$ Knoten \\
$m$ \>Kantenanzahl \\
$n$ \>Knotenanzahl \\
$N(v)$, $N_i(v)$ \>1-Umgebung bzw. $i$-Umgebung des Knotens $v$ \\
$s, u, u_i, v, v_i$ \>Knoten \\
$P_n$ \>Pfadgraph mit $n$ Knoten \\
$U$ \>Obere Schranke für die Brennzahl in der ILP-PROP-Formulierung\\
$V$ \>Knotenmenge \\
$w$ \>Kantengewicht \\
$Z(v)$ \>Zusammenhangskomponente des Knotens $v$ \\
\end{tabbing}
\subsection*{Zufallsvariablen, Verteilungen und Statistik}
\begin{tabbing}
xxxxxxxxxxxxxxxxxx\=xxxxxxxxxxxxxxxxxxxxxxxxxxxxxxxxxxxxx\kill
$X$, $Y$, $X_i$, $Y_i$ \>Zufallsvariablen \\
$\alpha$ \>Signifikanzniveau \\
$t_{\mathdutch{p}, \mathdutch{n}}$ \>$\mathdutch{p}$-Quantil der $t$-Verteilung mit $\mathdutch{n}$ Freiheitsgraden \\
$C$ \>Berry-Esseen-Konstante \\
$\chi^2_{\mathdutch{p}, n}$ \>$\mathdutch{p}$-Quantil der $\chi^2$-Verteilung mit $\mathdutch{n}$ Freiheitsgraden \\
$\E{X}$ \>Erwartungswert von $X$ \\
$\mu$, $\sigma$ \>Parameter der Normalverteilung \\
$\mu_X$, $\sigma_X$ \>Parameter der Normalverteilung von $X$ \\
$\gamma$ \>Schiefe \\
$\epsilon$ \>Fehlerterm \\
$f$ \>\pdf einer Normalverteilung \\
$F$ \>Verteilungsfunktion einer Normalverteilung \\
$f_X$ \>\pdf von $X$ \\
$F_X$ \>\cdf von $X$ \\
\end{tabbing}



\newpage
\cleardoublepage
\onehalfspacing
% Dateien
\section{Einleitung}
\pagenumbering{arabic}
Das Graph-Burning-Problem ist ein kombinatorisches Optimierungsproblem der theoretischen Informatik. 
Das Modell beschreibt einen zeitdiskreten Prozess, bei dem in jedem Schritt ein neuer Brandherd gelegt wird, während die aktive Feuerfront gleichzeitig auf die Nachbarknoten im Graphen übergreift. Das Ziel besteht darin, alle Knoten des Graphen in einer minimalen Anzahl von Zeitschritten zu verbrennen \cite{NAZERI2023}.

Das Problem ist für viele Anwendungen relevant. Es kann beispielsweise eingesetzt werden, um eine epidemiologische Ausbreitung \cite{Rosh16, GLM21}, einen Virusbefall in Computernetzwerken \cite{Rosh16}, die Verbreitung von Informationen in sozialen Netzwerken \cite{bjr14, Rosh16, Gau22, NAZERI2023, Ashok25} oder die Alarmierung der Bevölkerung \cite{Si19, SLP19, GLM21} im Zeitablauf zu modellieren. 
Ziel dieser Arbeit ist es, den gewählten Ansatz zu implementieren und unter Verwendung von verschiedenen Zentralitätsmaßen und Penalty-Funktionen empirisch zu untersuchen. 

Nach der Einleitung werden in Kapitel~\ref{cap:theo} die theoretischen Grundlagen vorgestellt. 
In Kapitel~\ref{cap:nazeri} wird dann der zu optimierende Ansatz von Nazeri et al. \cite{NAZERI2023} beschrieben. 
In Kapitel~\ref{cap:impl} folgt die Erläuterung der eigenen Implementierung. 
In Kapitel~\ref{cap:num} werden die numerischen Ergebnisse dargestellt. 
Die Arbeit wird in Kapitel~\ref{cap:conc} mit einem Fazit abgeschlossen. 

%%%%%%%%%%%%%%%%%
% Penalty-Verfahren
%%%%%%%%%%%%%%%%%
\section{Theoretische Grundlagen}\label{cap:theo}
We consider the Poisson equation on the unit square $\Omega = (0,1)^2$ with homogeneous Dirichlet boundary conditions:
\begin{eqnarray}
    %\begin{cases} 
        -\Delta u(x, y) = f & \text{for } (x, y)\in\Omega, \\ 
        u(x, y) = 0 & \text{for } (x, y)\in\partial\Omega. 
    %\end{cases}
\end{eqnarray}
We discretize the unit square with a regular grid with step size $h = 1/(N + 1)$. 
Thereby, we consider three stencils for $-\Delta u$: 

$S_5 = \frac{1}{h^2} \begin{bmatrix} 0 & -1 & 0 \\ -1 & 4 & -1 \\ 0 & -1 & 0 \end{bmatrix}$, 
$S_9 = \frac{1}{6h^2} \begin{bmatrix} -1 & -4 & -1 \\ -4 & 20 & -4 \\ -1 & -4 & -1 \end{bmatrix}$, and 
$S_5^{\text{diag}} = \frac{1}{2h^2} \begin{bmatrix} -1 & 0 & -1 \\ 0 & 4 & 0 \\ -1 & 0 & -1 \end{bmatrix}$. 

To compute the numerical solution, we apply the Jacobi method, which is notoriously known for its slow convergence. 
For the CPU implementation, we experiment with different block sizes, memory footprints, and core counts. 
For the GPU implementation, we \textbf{plan to} investigate different thread block dimensions, shared memory configurations, and memory access patterns.

As a stopping criterion, we use a maximum iteration count $k_{\max}$ as well as the $L_2$ norm of the Jacobi step size, $\|U - V\|_2$. 
The initial guess is generated randomly, and the exact same initial vector is used across all experiments to ensure comparability.

We test our implementation with different grid sizes and stencils, and measure the time for $k$ sweeps on the CPU and on the GPU, including the transfers.


\begin{align} \min \quad & \sum_{j \in [1, U]} x_{j}+ 1, \label{obj:ILP-PROP} \\ \text{so dass} \quad & x_{j} \geq 1 - b_{i,j} \qquad \forall v_i \in V, ~ j = 1,2,\ldots,U \label{res:ILP-PROP-xj}\\ & b_{i,j} \leq s_{i,j} + \sum_{v_l \in N(v_i)} b_{l,j-1} \qquad \forall v_i \in V, ~ j=1,2,\ldots,U \label{res:ILP-PROP-bij} \\ & \sum_{v_i \in V} s_{i,j} = 1 \qquad j=1,2,\ldots,U \label{res:ILP-PROP-sij}\\ 
& s_{i,j}, ~ b_{i,j}, ~ x_{j} \in \{0, 1\} \qquad \forall v_i \in V, ~ j=1,2,\ldots,U \label{res:ILP-PROP-bin}\\ & b_{i,0} = 0 \qquad \forall v_i \in V \label{res:ILP-PROP-bi0}
\end{align}

Das Beispiel von Goldberg \cite[11]{G89} (siehe Abbildung~\ref{fig:RWS}) verdeutlich das Prinzip der Roulette-Wheel-Selektion. 
Den vier verfügbaren Kandidaten werden Wahrscheinlichkeiten zugeordnet. 
Der ''beste'' Kandidat hat eine Wahrscheinlichkeit von 49.2\% gezogen zu werden, während der ''schwächste'' Kandidat nur mit einer Wahrscheinlichkeit von 5.5\% ausgewählt wird. 
\begin{table}[htb]
    \centering
    \small
    \begin{tabular}{rrrr}
        \hline
        Nr. & Kodierung & Fitness & \% \\
        \hline
        1 & 01101 & 169 & 14.4 \\
        2 & 11000 & 576 & 49.2 \\
        3 & 01000 & 64 & 5.5 \\
        4 & 10011 & 361 & 30.9 \\
        \hline
        Summe & & 1170 & 100.0 \\
        \hline
    \end{tabular}
\caption[Beispielhafte Fitnesswerte für Roulette-Wheel-Selektion]{Beispielhafte Fitnesswerte für Roulette-Wheel-Selektion \cite[11]{G89}}
\end{table}

Die Gesamtpopulation des implementierten genetischen Algorithmus besteht aus drei Teilen, die wie folgt aufgebaut sind:
\begin{table}[htb]
    \centering
    \small
    \begin{tabular}{lrr}
        \hline
        Populationsteil & Anzahl der Individuen & \% \\
        \hline
        Eliten & 60 & 10.0 \\
        Crossover & 450 & 75.0 \\
        Immigranten & 90 & 15.0 \\
        \hline
        Summe & 600 & 100.0 \\
        \hline
    \end{tabular}
\caption{Zusammensetzung der Gesamtpopulation des genetischen Algorithmus}
\end{table}

\subsection{Normalverteilung}
Die Normalverteilung ist eine zweiparametrige stetige Wahrscheinlichkeitsverteilung über der Menge der reellen Zahlen. Die Normalverteilung kann viele praktische Sachverhalte sehr gut approximieren. Dem zentralen Grenzwertsatz zufolge strebt die Summe von $\mathdutch{n}$ unabhängigen identisch verteilten Zufallsvariablen mit endlicher Varianz, unter schwachen Annahmen, mit $\mathdutch{n}\to\infty$ gegen eine Normalverteilung, unabhängig von der ursprünglichen Verteilung. Infolge dessen kann man Zufallsvariablen als annähernd normalverteilt ansehen, wenn sie durch Überlagerung von vielen unabhängigen kleinen Einflüssen entstehen \cite[309-310]{hen18}. 

\begin{Definition}[Normalverteilung]
Die Normalverteilung mit den Parametern $\mu \in \mathbb{R}$ und $0 < \sigma^2 < \infty$ hat für alle $x \in \mathbb{R}$ die Wahrscheinlichkeitsdichte
\begin{align}
	f(x) = \frac {1}{\sigma\sqrt{2\pi}} e^{-\frac {1}{2} \left(\frac{x-\mu}{\sigma}\right)^2}.
\end{align}		
Die \cdf der Normalverteilung ist durch
\begin{align}
	F(x) = \frac{1}{\sigma\sqrt{2\pi}} \int_{-\infty}^x e^{-\frac{1}{2} \left(\frac{t-\mu}{\sigma}\right)^2} \mathrm dt
\end{align}
gegeben. 
Eine Normalverteilung mit den Parametern $\mu = 0$ und $\sigma^2 = 1$ heißt Standardnormalverteilung. 
\end{Definition}
Zur Abkürzung wird oft $X \sim \mathcal N\left(\mu, \sigma^2\right)$ geschrieben. $\mu$ entspricht dabei dem Erwartungswert und $\sigma^2$ der Varianz der Zufallsgröße $X$.

Jede beliebige Normalverteilung kann standardisiert werden. Es sei $X \sim \mathcal N (\mu, \sigma^2)$. Die Transformation $Z = \frac{X - \mu}{\sigma}$ führt zu $Z \sim \mathcal N (0, 1)$. Die Rücktransformation der Ergebnisse mit $x = \mu + \sigma z$ ermöglicht die Berechnungen für beliebige Normalverteilungen \cite[309-310]{hen18}:  
\begin{align}
	\Phi(z) = \frac{1}{\sqrt{2\pi}} \int\limits_{-\infty}^{z} e^{-\frac{1}{2} t^2} \mathrm dt = \frac{1}{\sqrt{2\pi}} \int\limits_{-\infty}^{(x-\mu)/\sigma} e^{-\frac{1}{2} t^2} \mathrm dt = \Phi \left(\frac{x-\mu}{\sigma}\right).
\end{align}
Obwohl keine elementare Stammfunktion der \pdf der Normalverteilung existiert, sind die Quantile der Standardnormalverteilung $\mathcal N (0,1)$ tabelliert \cite[Section 1.3.6.7.1]{nist}. Computer-Algebra-Systeme beherrschen nicht nur diese tabellierten Quantile, sondern auch zahlreiche effiziente Algorithmen zur Bestimmung von numerischen Näherungswerten, um den gesuchten Wert mit Hilfe von schnell berechenbaren Approximationen zu ermitteln \cite[931-933]{as72}. 

Die \cdf der Standardnormalverteilung kann auf die (Gaußsche) Fehlerfunktion zurückgeführt werden \cite[298, 7.1.22]{as72}:  
\begin{align}
	P(Z \leq z) &= \Phi (z) = \frac{1}{2} \left[ 1+ \erf \left( \frac{z}{\sqrt{2}} \right) \right]. 
\end{align}

%%%%%%%%%%%%%%%%%
% Lösung von Nazeri et al. 
%%%%%%%%%%%%%%%%%
\section{CPU Results}
\lstset{%
    language         = C++,
    inputencoding    = utf8,
    extendedchars=true,
    literate=%
    {ő}{{{\H{o}}}}1
        {Ő}{{{\H{O}}}}1
        {ö}{{\"o}}1 {ä}{{\"a}}1 {ü}{{\"u}}1 
        {Ö}{{\"O}}1 {Ä}{{\"A}}1 {Ü}{{\"U}}1 
        {ß}{{\ss}}1 
         {ű}{{\H{u}}}1 {Ű}{{\H{U}}}1 % Hungarian fix
        {á}{{\'a}}1 {é}{{\'e}}1 {í}{{\'i}}1 {ó}{{\'o}}1 {ú}{{\'u}}1,
    basicstyle       = \scriptsize\ttfamily,
    keywordstyle     = \bfseries\color{blue},
    stringstyle      = \color{magenta},
    commentstyle     = \color{ForestGreen},
    showstringspaces = false,
}
\lstset{ 
  basicstyle=\scriptsize\ttfamily\linespread{0.7}\selectfont,        % the size of the fonts that are used for the code
  breakatwhitespace=false,         % sets if automatic breaks should only happen at whitespace
  breaklines=true,                 % sets automatic line breaking
  captionpos=b,                    % sets the caption-position to bottom
  %commentstyle=\itshape\color{red},    % comment style
 % deletekeywords={...},            % if you want to delete keywords from the given language
%  escapeinside={\%*}{*)},          % if you want to add LaTeX within your code
  extendedchars=true,              % lets you use non-ASCII characters; for 8-bits encodings only, does not work with UTF-8
  %firstnumber=1000,                % start line enumeration with line 1000
  frame=single,	                   % adds a frame around the code
  keepspaces=true,                 % keeps spaces in text, useful for keeping indentation of code (possibly needs columns=flexible)
 % keywordstyle=\color{blue},       % keyword style
  %language=Octave,                 % the language of the code
  %morekeywords={*,...},            % if you want to add more keywords to the set
  numbers=none,                    % where to put the line-numbers; possible values are (none, left, right)
  numbersep=5pt,                   % how far the line-numbers are from the code
  %numberstyle=\tiny\color{mygray}, % the style that is used for the line-numbers
  rulecolor=\color{black},         % if not set, the frame-color may be changed on line-breaks within not-black text (e.g. comments (green here))
  showspaces=false,                % show spaces everywhere adding particular underscores; it overrides 'showstringspaces'
  showstringspaces=false,          % underline spaces within strings only
  showtabs=false,                  % show tabs within strings adding particular underscores
  stepnumber=5,                    % the step between two line-numbers. If it's 1, each line will be numbered
 % stringstyle=\color{mymauve},     % string literal style
  tabsize=1 ,	                   % sets default tabsize to 2 spaces
  title=\lstname                   % show the filename of files included with \lstinputlisting; also try caption instead of title
%  caption={~/Downloads/Kunoth-21/julia/Blatt07/qrsvd.jl}
}
%%%%%%%%%%%%%%%%%
% \subsection{Evaluierung}
%%%%%%%%%%%%%%%%%
\subsection{Evaluierung}
Ein Knoten $v_i \in V$ gilt als verbrannt, wenn er zu irgendeinem Zeitpunkt des Brennprozesses entzündet wurde. 
Entweder wird er direkt angezündet oder er fängt Feuer durch einen direkten Nachbarn, der bereits brennt. 
Am Ende des Prozesses wird der Abstand von Knoten $v_i$ zum nächsten verbrannten Knoten als Restabstand $d_i$ bezeichnet. 
Für jeden verbrannten Knoten gilt $d_i = 0$. 
Eine valide Brennsequenz zeichnet sich dadurch aus, dass alle Knoten $v_i \in V$ vollständig verbrannt werden. 
Der Preis einer Brennsequenz $S = (v_1, v_2, \ldots, v_b)$ wird als die Summe der quadrierten Restabstände aller Knoten definiert \cite[6]{NAZERI2023}: 
\begin{align}
\label{eq:Cost}
p(S) = \sum_{v_i \in V} d_i^2,
\end{align}
wobei $d_i$ wie folgt berechnet werden kann:   
\begin{align}
\label{eq:d-i}
d_i = 
\begin{cases} 
  0 & \text{falls Knoten } v_i \text{ verbrannt ist,} \\
  \min_{1 \leq j \leq b} \, \{\text{dist}(v_i, v_j) - (b - j) \} & \text{sonst.}
\end{cases}
\end{align} 

Aufgrund der spezifischen Eigenschaften des GBP werden dafür neue Crossover- und Mutationsoperatoren eingeführt, die auf folgenden Eigenschaften basieren \cite[7]{NAZERI2023}: 
\begin{itemize}
\item Die Reihenfolge der Brandherde wird durch diese Operationen nicht verändert. 
\item Die Nachbarschaften der jeweiligen Brandherde werden durch diese Operationen exploriert.
\item Die Exploration der Nachbarschaft eines Knotens soll von seiner Zentralität abhängen.
\end{itemize}


%%%%%%%%%%%%%%%%%
% Implementierung
%%%%%%%%%%%%%%%%%
\section{GPU Results}
Der Algorithmus wird in C++ unter Verwendung von NetworKit 11.2 implementiert. 
Als Ausgangsbasis dient der bereitgestellte Code aus \cite{NAZERI2023}. 
Zusätzlich wird Python verwendet, um zufällige Graphen zu generieren, Ergebnisse zu visualisieren sowie nicht-tabellierte Quantile statistischer Verteilungen zu berechnen. 
Um das Programm in verschiedenen Konfigurationen und auf unterschiedlichen Datensätzen auszuführen, wird ein Bash-Skript erstellt. 
Der vollständige Quellcode ist im GitHub-Repository \url{https://github.com/serin2001/GBP-GA} verfügbar. 

%\subsection{Voraussetzungen}
Die Hardwareausstattung und die installierte Software, die für die Entwicklung und Durchführung der Experimente verwendet wurden, sind nachfolgend zusammengefasst:
%\input{Chap04_tab_voraussetz.tex}

%%%%%%%%%%%%%%%%%
% Probleminstanzen
%%%%%%%%%%%%%%%%%
\subsection{Instanzen}
Der Algorithmus wird anhand von mehreren Datensätzen evaluiert. 
Alle Graphen werden als ungerichtet und ungewichtet behandelt. Im Rahmen des Preprocessings werden gerichtete oder gewichtete Graphen transformiert. Gerichtete Kanten werden als ungerichtet eingelesen. Anschließend werden Kantengewichte, Mehrfachkanten und Schleifen entfernt. Bei den mit $^*$ markierten Datensätzen ergeben sich dadurch die angegebenen Kantenzahlen. Alle verwendeten Graphen mit Ausnahme von \texttt{cite} sind zusammenhängend. 

Die nachfolgenden 13 Graphen stammen aus öffentlich zugänglichen Datenbanken \cite{snapnets, nr}. Sie dienen als Benchmark-Instanzen für den Vergleich der Güte verschiedener Lösungsverfahren. Tabelle~\ref{tab:Chap04_tab_benchmarks} vermittelt einen Überblick über diese Testinstanzen. Die Spalte Dateiname gibt den entsprechenden Suchbegriff für die jeweilige Datenbank und ggf. die Bezeichnung innerhalb der Archivdatei an. 

%\input{Chap04_tab_benchmarks}

\lstset{%
    language         = C++,
    inputencoding    = utf8,
    extendedchars=true,
    literate=%
    {ő}{{{\H{o}}}}1
        {Ő}{{{\H{O}}}}1
        {ö}{{\"o}}1 {ä}{{\"a}}1 {ü}{{\"u}}1 
        {Ö}{{\"O}}1 {Ä}{{\"A}}1 {Ü}{{\"U}}1 
        {ß}{{\ss}}1 
         {ű}{{\H{u}}}1 {Ű}{{\H{U}}}1 % Hungarian fix
        {á}{{\'a}}1 {é}{{\'e}}1 {í}{{\'i}}1 {ó}{{\'o}}1 {ú}{{\'u}}1,
    basicstyle       = \scriptsize\ttfamily,
    keywordstyle     = \bfseries\color{blue},
    stringstyle      = \color{magenta},
    commentstyle     = \color{ForestGreen},
    showstringspaces = false,
}
\lstset{ 
  basicstyle=\scriptsize\ttfamily,        % the size of the fonts that are used for the code
  breakatwhitespace=false,         % sets if automatic breaks should only happen at whitespace
  breaklines=true,                 % sets automatic line breaking
  captionpos=b,                    % sets the caption-position to bottom
  extendedchars=true,              % lets you use non-ASCII characters; for 8-bits encodings only, does not work with UTF-8
  frame=single,	                   % adds a frame around the code
  keepspaces=true,                 % keeps spaces in text, useful for keeping indentation of code (possibly needs columns=flexible)
  numbers=none,                    % where to put the line-numbers; possible values are (none, left, right)
  rulecolor=\color{black},         % if not set, the frame-color may be changed on line-breaks within not-black text (e.g. comments (green here))
  showspaces=false,                % show spaces everywhere adding particular underscores; it overrides 'showstringspaces'
  showstringspaces=false,          % underline spaces within strings only
  showtabs=false,                  % show tabs within strings adding particular underscores
  stepnumber=5,                    % the step between two line-numbers. If it's 1, each line will be numbered
  tabsize=1 ,	                   % sets default tabsize to 2 spaces
  title=\lstname                   
}

Der Algorithmus wurde auch auf zufällig generierten Graphen getestet. Es wurden 100 Erd\H{o}s--R\'enyi-Graphen mit $1000$ Knoten und 6000 Kanten sowie 100 Barab\'asi--Albert-Graphen mit $1000$ Knoten und 3000 Kanten generiert. Diese Anzahl wurde im Hinblick auf Konfidenzintervalle gewählt. Während alle generierten Erd\H{o}s--R\'enyi-Graphen eine Brennzahl $b(G) = 5$ aufweisen, zeigen die Barab\'asi--Albert-Graphen eine Streuung. 
%\input{Chap04_fig_BA_ER}
%Die beiden Zufallsgraphentypen werden in Abbildungen~\ref{fig:BA_plot} und \ref{fig:ER_plot} visualisiert. 

Zu Testzwecken während der Entwicklung wurden zudem die Pfadgraphen $P_{64}$ und $P_{81}$ hinzugefügt. Für die optimale Länge der Brennsequenz gilt: $b(P_{n})=\lceil \sqrt{n}\rceil$. 
Für solche Pfadgraphen sind effiziente Lösungsverfahren verfügbar.  
%%%%%%%%%%%%%%%%%
% Darstellung
%%%%%%%%%%%%%%%%%
\subsubsection{Dateiformat}
Nachfolgend wird ein Ausschnitt von \texttt{P64.csv} erklärt. Der Graph besitzt 64 Knoten und 63 Kanten. Die Nummerierung der Knoten beginnt bei 1. Ungewichtet bedeutet, dass ab der dritten Zeile keine Kantengewichte angegeben werden, die das Programm verwerfen muss. Der entsprechende Graph wird in Abbildung~\ref{fig:pfadgraph_pn} dargestellt. 

%\lstinputlisting[language=C, firstline=1, lastline=5, caption={Aufbau einer Testinstanz anhand von $P_{64}$}, 
%frame=none, %resetmargins=true, 
%linewidth=1.0\linewidth]{./Data/P64.csv}
\begin{lstlisting}[
    language={}, 
    frame=none,
    caption={Beschreibung einer Instanz anhand von \texttt{P64.csv} (Ausschnitt)}, 
    label={lst:p64_aufbau}
]
P64
64 63 1 0
1 2
2 3
3 4
\end{lstlisting}
%\input{Chap04_fig_pfadgraph}
Wäre der Graph gewichtet, so müssten für jede Kante drei durch Leerzeichen getrennte Einträge angegeben werden: die beiden Endknoten und das jeweilige Kantengewicht. Auf eine separate Darstellung des gewichteten Dateiformats wird verzichtet. 

%%%%%%%%%%%%%%%%%
% Ausführung
%%%%%%%%%%%%%%%%%
\subsection{Ausführung des Programms}
Das Programm kann auf zwei Arten ausgeführt werden: Es kann entweder direkt oder über ein mitgeliefertes Bash-Skript aufgerufen werden. In beiden Fällen können optionale Parameter übergeben werden, um eine selektive Ausführung zu starten. 

Beispielhaft wird die Testinstanz \texttt{P64.csv} auf eine Brennsequenz der Länge 8 untersucht: 

\texttt{./CBAG P64.csv 8} 

Die Berechnungen werden für alle Graphen in Tabelle~\ref{tab:Chap04_tab_benchmarks} 20-mal wiederholt. 
Da eine mehrfache Berechnung der insgesamt 200 stochastischen Graphen zu einer Verlängerung der Gesamtlaufzeit führt, werden diese Instanzen nur einmal berechnet. 
Falls eine Wiederholung dennoch gewünscht ist, muss \texttt{REPEAT=1} in Zeile 128 von \texttt{run.sh} auskommentiert werden.
%%%%%%%%%%%%%%%%%
% Optionale Parameter
%%%%%%%%%%%%%%%%%
\subsubsection{Optionale Parameter}
Um eine gezielte Auswahl von Zentralitätsmaßen und/oder Penalty-Funktionen zu treffen sowie den Umfang der Log-Ausgabe zu steuern, können Parameter \texttt{-c}, \texttt{-p}, \texttt{-r} und \texttt{-v} übergeben werden. 

\texttt{./run.sh -p SQRT QUADRATIC -c DEGREE CLOSENESS -v 65 -r 10}

Die nachfolgende Tabelle \ref{tab:options-c} listet alle implementierten Zentralitätsmaße auf, die als optionale Parameter mit \texttt{-c} beliebig kombiniert übergeben werden können. Das Weglassen von \texttt{-c} führt zur Benutzung aller aufgeführten Zentralitätsmaße.
%\input{Chap04_tab_c_option}

In der Tabelle bezeichnet $\mathdutch{u}=\{v_1, \ldots, v_{k}\}$ die Menge der unverbrannten Knoten, wobei $k$ deren Anzahl zum Zeitpunkt der Bewertung angibt, und $d_i$ den kürzesten erreichbaren Abstand des Knotens $v_i \in \mathdutch{u}$ zur aktiven Feuerfront. 
Alle %sonstigen 
Penalty-Funktionen $p$ haben die Gestalt 
\begin{equation}
%\begin{split}
p \colon \mathbb N_{0}^{|\mathdutch{u}|} \to \mathbb R_{\geq 0}, \quad \mathdutch{u} \mapsto \begin{cases} 
\xi \displaystyle\sum_{v_i \in \mathdutch{u}} d_i^j \text{ laut Tabelle \ref{tab:Chap04_tab_p_option}, } & \text{falls } |\mathdutch{u}| < \texttt{threshold} \\
c\xi \,|\mathdutch{u}|,  & \text{falls } |\mathdutch{u}| \geq \texttt{threshold}.
\end{cases}
%\end{split}
\label{eq:penalty}
\end{equation}
Die Konstante $c$ entspricht dabei dem Wert \texttt{INF} im Programm und wird in Abhängigkeit von $j$ und \texttt{threshold} mindestens so hoch gewählt wie die Summe der ersten $|\mathdutch{u}|$ natürlichen Zahlen in der $j$-ten Potenz 
\begin{equation}
\sum_{1 \leq i \leq |\mathdutch{u}|} i^j 
\end{equation}
ausgewertet für $|\mathdutch{u}| = \texttt{threshold}$. 
Diese Summe kann mit einer geschlossenen Formel (siehe~\cite{wolfram}) berechnet werden. 
Diese Wahl soll eine konsistente Penalty-Funktion sicherstellen. Eine Penalty-Funktion muss hierbei nicht notwendigerweise stetig sein; sie muss lediglich schlechtere Chromosomen stärker bestrafen. Mit der Definition (\ref{eq:penalty}) können auch weit von einer zulässigen Lösung entfernte Chromosomen sinnvoll gegeneinander bewertet werden. 

%%%%%%%%%%%%%%%%%
% Programmstruktur
%%%%%%%%%%%%%%%%%
\subsection{Klassenstruktur des Programms}
%%%%%%%%%%%%%%%%%
% main.cpp
%%%%%%%%%%%%%%%%%
%\subsubsection{Makefile und main.cpp}
C++ ist eine Programmiersprache, die kompiliert werden muss. 
Dies erfolgt über \texttt{Makefile} (siehe Listing~\ref{lst:Makefile} im Anhang). 
Die Erstellung dieser Datei erfolgte mithilfe von \url{google.ai}. 
Die Kompilierung muss nach jeder Programmänderung mit \texttt{make} angestoßen werden.  

%%%%%%%%%%%%%%%%%
% CentralityComputer.cpp
%%%%%%%%%%%%%%%%%
\subsubsection{CentralityComputer}
Die Klasse \texttt{CentralityComputer} (siehe Listings~\ref{lst:CentralityComputer.cpp} und \ref{lst:CentralityComputer.hpp} im Anhang) berechnet die Zentralitätswerte für jeden Knoten $v \in V$ und skaliert diese auf das Intervall [0,\,1], damit Knoten aus verschiedenen Zusammenhangskomponenten vergleichbar oft ausgewählt werden. %Zusätzlich wendet sie eine Sigmoidfunktion auf die berechneten Rohwerte an. 
Die Klasse operiert mit einem Pointer, der auf das \texttt{GraphAnalyzer}-Objekt zeigt. 
%Dabei wird lediglich dessen Speicheradresse im RAM übergeben.

Für die Bestimmung der Zentralitätswerte wird auf NetworKit zurückgegriffen. 
Die Auswahl der implementierten Zentralitätsmaße (siehe Tabelle~\ref{tab:options-c}) richtet sich nach den folgenden Kriterien: 

\begin{itemize}
\item \texttt{BETWEENNESS}, \texttt{DEGREE}, \texttt{CLOSENESS} und \texttt{EIGENVECTOR} zählen zu den häufigsten Zentralitätsmaßen in der Literatur \cite{Gau22, NAZERI2023, Ko05} sowie \cite[Kap. 7]{New18} und \cite[Kap. 6]{J03}.
\item \texttt{BETWEENNESS} und \texttt{EIGENVECTOR} werden in vergleichbaren Studien \cite{Gau22, NAZERI2023} als zentrale Maße eingesetzt.
\item Die Zentralitätsmaße \texttt{PAGERANK}, \texttt{EIGENVECTOR} und \texttt{KATZ} sind eng verwandt. 
\item \texttt{CLOSENESS} und \texttt{HARMONIC} sind eng verwandt. 
\item \texttt{KADABRA} und \texttt{BETWEENNESS} sind eng verwandt. 
\item \texttt{DEGREE} lässt sich in puncto Rechenaufwand am günstigsten bestimmen. 
\end{itemize}

Dabei wird von den verwandten Zentralitätsmaßen eine vergleichbare Leistung erwartet. 
Als Kontrollmaß dafür, ob die Zentralität überhaupt einen Effekt auf den genetischen Algorithmus hat, wird \texttt{RANDOM} verwendet. Dabei werden allen Knoten $v \in V$ gleichverteilte Zufallswerte aus dem Intervall [0,~1] zugewiesen. 

Die Verwendung der NetworKit-Bibliothek bietet zudem den Vorteil, dass zusätzliche Zentralitätsmaße einfach implementiert werden können. 
Exemplarisch soll das Katz-Zentralitätsmaß in das Programm integriert werden.   
Dafür müssen lediglich 

{\small{\texttt{\#include <networkit/centrality/KatzCentrality.hpp>}}}
 
und der entspechende Codeblock (siehe Listing~\ref{lst:computeCentralityScores}) in die Methode \texttt{computeCentralityScores} eingefügt werden. Zudem muss die neue Abkürzung an je zwei Stellen in den Dateien \texttt{main.cpp} und \texttt{CentralityComputer.hpp} in den Listen der erlaubten Werte ergänzt werden. Die optionale Skalierung der Werte mit einem konstanten Faktor kann dazu beitragen, eine Division durch sehr kleine Werte zu vermeiden. 

%\input{Chap04_lst_computeCentralityScores}

Im Anschluss wird die Zentralität $c_i$ jedes Knotens $v_i \in V$ durch das Maximum $\max_{v_j \in Z(v_i)} c_j$ seiner Zusammenhangskomponente $Z(v_i)$ geteilt. 
In den trivialen Fällen $Z(v) = \{v\}$ oder $Z(v) = \{v, u\}$ wird allen Knoten aus $Z(v)$ eine normierte Zentralität von 1 zugewiesen. 
%Dies trägt dazu bei, dass die Chromosomen Brandherde aus allen Zusammenhangskomponenten enthalten. 

Eine weitere wichtige Funktion der Klasse \texttt{CentralityComputer} ist das Bereitstellen eines statischen Zufallszahlengenerators des Typs \texttt{mt19937\_64} \cite{mt19937}, der hier bei der \texttt{RANDOM}-Strategie und in der Klasse \texttt{GA} zum Einsatz kommt. 
Im Gegensatz zur Lösung von Nazeri et al. \cite{NAZERI2023} wird nur ein 64"=Bit"=Zufallszahlengenerator erzeugt, der je nach Bedarf ganzzahlige oder reellwertige Zufallszahlen in dem gewünschten Intervall liefern kann.  
 


%%%%%%%%%%%%%%%%%
% Numerische Ergebnisse
%%%%%%%%%%%%%%%%%
\section{Results}
All experiments were conducted on an \textbf{AMD Ryzen 7 5700X} CPU (8 cores, 32\,MB L3 cache, measured STREAM triad bandwidth of $36.2\,\text{GB/s}$) and an \textbf{AMD Radeon RX 9070} GPU (measured STREAM triad bandwidth of $682.45\,\text{GB/s}$). 

Across both architectures, all three stencils ($S_5$, $S_9$, and $S_5^{\text{diag}}$) achieve nearly identical throughput, confirming that the Jacobi kernels are strictly \emph{memory-bound}. On the large $2^{24}$-point grid, the spatially blocked CPU implementation sustains $\approx 1500\,\text{MUpdates/s}$, effectively saturating the Ryzen 7 5700X's main memory bandwidth. Without spatial blocking, CPU performance degrades by up to $40\%$ (dropping to $\approx 891\,\text{MUpdates/s}$) once the row length reaches $N_x \ge 65\,536$. At this size, three consecutive grid rows exceed the L1/L2 cache capacity, leading to cache eviction and redundant DRAM accesses. Spatial blocking eliminates this layer-condition bottleneck entirely.

On the Radeon RX 9070, the pure compute kernel maintains a stable throughput of $\approx 13.2$ GUpdates/s across all grid aspect ratios---an $\approx 8.8\times$ speedup over the blocked CPU execution, driven by the GPU's $18.85\times$ higher memory bandwidth. While PCIe host-to-device data transfers noticeably reduce the effective throughput during short 50-sweep runs ($6.0\text{--}10.9\,\text{GUpdates/s}$), this transfer overhead becomes negligible ($< 2\%$) over $1000$ sweeps. On the smaller $512 \times 512$ grid ($2\,\text{MB}$ per double-precision array), the working set fits entirely within the Ryzen 7 5700X's $32\,\text{MB}$ L3 cache, boosting CPU throughput to $\approx 5.0\text{--}5.4\,\text{GUpdates/s}$, whereas the GPU reaches $\approx 10.5\,\text{GUpdates/s}$ due to kernel launch latencies.

Finally, comparing numerical convergence on the $512 \times 512$ grid after $1000$ sweeps highlights the efficiency of the compact 9-point stencil $S_9$: because its additional arithmetic operations are hidden by the memory-bandwidth limit, $S_9$ incurs no runtime penalty while reducing the $L_2$ residual by more than two orders of magnitude compared to $S_5$ ($0.099$ vs.\ $12.76$).

\subsection*{Performance and Convergence Analysis}

All experiments were conducted on an \textbf{AMD Ryzen~7~5700X} CPU (8 cores, $32\,\text{MB}$ L3 cache, measured STREAM triad bandwidth of $37.4\,\text{GB/s}$) and an \textbf{AMD Radeon~RX~9070} GPU ($8\,\text{MB}$ L2 cache, $64\,\text{MB}$ Infinity Cache, measured STREAM triad bandwidth of $684.28\,\text{GB/s}$).

On the large $2^{24}$-point grid ($128\,\text{MB}$ per double-precision array), all three stencils ($S_5$, $S_9$, and $S_5^{\text{diag}}$) achieve nearly identical throughput across both architectures, confirming that the kernels are strictly \emph{memory-bound} (Table~\ref{tab:aspect_ratio}). The spatially blocked CPU version sustains $\approx 1.50\text{--}1.54\,\text{GUpdates/s}$, saturating the $37.4\,\text{GB/s}$ main memory bandwidth at the theoretical minimum of $24\,\text{Bytes/update}$. Without spatial blocking, CPU performance drops by up to $41\%$ (to $\approx 0.89\,\text{GUpdates/s}$) for row lengths $N_x \ge 65\,536$, as three consecutive grid rows exceed the L1/L2 cache capacity and trigger cache evictions. Spatial blocking eliminates this bottleneck completely.

\begin{table}[htbp]
    \centering
    \caption{Throughput (in GUpdates/s) across varying grid aspect ratios ($N_x \times N_y = 2^{24}$) for the compact 9-point stencil $S_9$ over $50$ sweeps.}
    \label{tab:aspect_ratio}
    \small
    \begin{tabular}{rrcccc}
        \hline
        \textbf{$N_x$} & \textbf{$N_y$} & \textbf{CPU (unblocked)} & \textbf{CPU (blocked)} & \textbf{GPU (pure)} & \textbf{GPU (transfer)} \\
        \hline
        $256$         & $65\,536$ & $1.53$ & $1.54$ & $13.16$ & $11.21$ \\
        $4096$        & $4096$    & $1.51$ & $1.51$ & $13.14$ & $10.23$ \\
        $32\,768$     & $512$     & $1.32$ & $1.32$ & $13.13$ & $10.91$ \\
        $65\,536$     & $256$     & $1.10$ & $1.42$ & $13.12$ & $10.28$ \\
        $262\,144$    & $64$      & $0.89$ & $1.42$ & $13.22$ & $10.91$ \\
        $2\,097\,152$ & $8$       & $1.38$ & $1.42$ & $13.15$ & $9.63$  \\
        \hline
    \end{tabular}
\end{table}

On the Radeon~RX~9070, the pure compute kernel maintains a stable throughput of $\approx 13.1\text{--}13.4\,\text{GUpdates/s}$ across all aspect ratios on the $2^{24}$ grid---an $\approx 8.8\times$ speedup over the blocked CPU implementation. While PCIe host-to-device transfers reduce effective throughput during short 50-sweep runs ($9.6\text{--}11.4\,\text{GUpdates/s}$), this overhead drops to $< 3.5\%$ over $1000$ sweeps (Table~\ref{tab:convergence_perf}).

\begin{table}[htbp]
    \centering
    \caption{Throughput (in GUpdates/s) and final $L_2$ residual norm after $1000$ Jacobi sweeps on the Ryzen~7~5700X (blocked) and Radeon~RX~9070.}
    \label{tab:convergence_perf}
    \small
    \begin{tabular}{llcccc}
        \hline
        \textbf{Grid Size} & \textbf{Stencil} & \textbf{CPU} & \textbf{GPU (pure)} & \textbf{GPU (total)} & \textbf{$L_2$ Residual} \\
        \hline
        $512 \times 512$   & $S_5$               & $4.93$ & $10.47$ & $10.27$ & $12.764$ \\
                           & $S_9$               & $5.25$ & $10.48$ & $10.36$ & $0.099$  \\
                           & $S_5^{\text{diag}}$ & $5.36$ & $10.70$ & $10.56$ & $6.914$  \\
        \hline
        $2048 \times 2048$ & $S_5$               & $2.52$ & $13.41$ & $12.95$ & $57.767$ \\
                           & $S_9$               & $2.54$ & $16.56$ & $16.39$ & $0.199$  \\
                           & $S_5^{\text{diag}}$ & $2.52$ & $22.18$ & $21.88$ & $29.237$ \\
        \hline
    \end{tabular}
\end{table}

Comparing the $512 \times 512$ ($2\,\text{MB}$/array) and $2048 \times 2048$ ($32\,\text{MB}$/array) grids over $1000$ sweeps reveals distinct cache and latency regimes:
\begin{itemize}
    \item \textbf{CPU Cache Hierarchy:} On the $512 \times 512$ grid, the working set fits entirely inside the Ryzen~7~5700X's $32\,\text{MB}$ L3 cache, yielding $\approx 4.9\text{--}5.4\,\text{GUpdates/s}$. At $2048 \times 2048$, the arrays partially spill into main memory, settling at an intermediate rate of $\approx 2.52\,\text{GUpdates/s}$.
    \item \textbf{GPU Launch Latency vs.\ Infinity Cache:} Despite fitting within the GPU's $8\,\text{MB}$ L2 and $64\,\text{MB}$ L3 caches, the $512 \times 512$ grid is limited to $\approx 10.5\,\text{GUpdates/s}$ by kernel launch overhead ($\approx 25\,\mu\text{s}$ per sweep) and low compute-unit occupancy ($262\text{k}$ threads). At $2048 \times 2048$ ($4.19\text{M}$ threads), launch latencies are amortized and the working set still benefits from the $64\,\text{MB}$ Infinity Cache, pushing throughput up to $13.4\text{--}22.2\,\text{GUpdates/s}$ (up to an $8.8\times$ GPU speedup).
    \item \textbf{Numerical Convergence:} Across both grid resolutions, the compact 9-point stencil $S_9$ incurs no runtime penalty over $S_5$ while reducing the final $L_2$ residual after $1000$ sweeps by more than two orders of magnitude ($0.199$ vs.\ $57.767$ at $2048 \times 2048$).
\end{itemize}


%%%%%%%%%%%%%%%%%
% Fazit
%%%%%%%%%%%%%%%%%
\section{Fazit}\label{cap:conc}
Die vorangegangenen Darstellungen verdeutlichen, wie problematisch die Frage nach der empirisch besten Strategie für einen stochastischen Algorithmus ist, wenn keine hinreichend große Stichprobe generiert werden kann. 

Zudem hat sich gezeigt, dass von einer parallelen Berechnung Abstand genommen werden sollte. 
Die parallel verarbeiteten Instanzen weisen tendenziell verschlechterte Gesamtlaufzeiten auf. 
Die Brennzahlen aus \cite{NAZERI2023} konnten für jede einzelne Instanz und jede Strategie bestätigt, aber in keinem Fall unterboten werden. 
Die entscheidende Erkenntnis bleibt die Zuverlässigkeit der implementierten Lösung: 
Das Programm kann für jede der 45 untersuchten Kombinationen aus Zentralitätsmaß und Penalty-Funktion zulässige Lösungen generieren. 
Aufgrund der exzessiv langen Laufzeiten des Programms kann jedoch keine ausreichend große Stichprobe für alle 45 Strategien gewonnen werden. 

\cleardoublepage
\appendix
\singlespacing
%\input{appendix}
\newpage
     	
%\printbibliography
%\addcontentsline{toc}{section}{Literatur}
%\newpage
%\cleardoublepage
%\input{Erklaerung}
%\addcontentsline{toc}{section}{Erklärung an Eides statt}
\end{document}
